%%%% utfpr-poster.tex, 2024/10/05
%%%% Copyright (C) 2018-2024 Luiz E. M. Lima (luizeduardomlima@gmail.com)
%%
%% This work may be distributed and/or modified under the conditions of the
%% LaTeX Project Public License, either version 1.3 of this license or (at your
%% option) any later version.
%% The latest version of this license is in
%%   https://www.latex-project.org/lppl.txt
%% and version 1.3 or later is part of all distributions of LaTeX version
%% 2005/12/01 or later.
%%
%% This work has the LPPL maintenance status `maintained'.
%%
%% The Current Maintainer of this work is Luiz E. M. Lima.
%%
%% This work consists of all files listed in readme.html.
%%
%% A brief description of this work is in readme.txt.

%% Detecção e aviso sobre comandos obsoletos
% \RequirePackage[l2tabu, orthodox]{nag}

%% Classe de documento
\documentclass[%% Opções (^ = padrão; > = para pacotes):
  t,%% Alinhamento vertical de quadros: b (base), c (centro) (^) ou t (topo)
  onlytextwidth,%% Colunas ocupando somente largura de texto
  compress,%% Redução de tamanho de barras de navegação
%   draft,%% Aparência de documento (>): draft ou final (^)
  xcolor = table,%% Utilização do pacote colortbl (>)
  english,%% Idioma secundário (penúltimo) (>)
  brazilian,%% Idioma primário (último) (>)
]{beamer}

%% Pacotes utilizados
\usepackage[%% Opções (^ = padrão; > = para beamerposter):
%   Layout    = Beamer,%% Leiaute (blocos): TCB (tcolorbox) (^) ou Beamer
%   Font      = Arial,%% Fonte principal: Arial, Calibri, Times ou LaTeX(*) (^)
%   LinkColor = ,%% Cor (Modelo HTML) hyperlink (azul escuro ^): nenhum modelo (cor texto)
%   Justify   = Off,%% Texto justificado: On (^) ou Off
%   CCLic     = BY-NC-ND,%% Licença CC: BY (^), BY-SA, BY-ND, BY-NC, BY-NC-SA ou BY-NC-ND
%   ABNTCit   = NSB,%% Citação ABNT: AAY (Nome, Ano) (^), NRB (1) ou NSB [1]
%   DOIIcon   = On,%% Ícone de DOI em referências: On ou Off (^)
%   URLIcon   = On,%% Ícone de URL em referências: On ou Off (^)
  size      = custom,%% Tamanho de papel (>): a0 (^), a1, a2, a3, a4 ou custom
  width     = 100,%% Largura de papel em centímetros (custom) (>)
  height    = 120,%% Altura de papel em centímetros (custom) (>)
  scale     = 1.5,%% Escala de fontes (>): 1.0 (^), 1.5, etc.
]{utfpr-poster}

%% Plano de fundo no pôster (arquivos de imagem em ./Backgrounds/)
\PosterBackGround{utfpr-b-bg-100-x-120}

%% Informações do documento (descomentar para alterar; * = opcional)
%%%% Título
\title{%
  Título do Pôster para Apresentação em Congresso,\\%
  Seminário ou Evento Técnico/Científico%
}
%%%% Subtítulo (*)
\subtitle{%
  Subtítulo do Pôster para Apresentação em Congresso,\\%
  Seminário ou Evento Técnico/Científico%
}
%%%% Assunto
\subject{Nome e/ou Sigla do Congresso, Seminário ou Evento Técnico/Científico}
%%%% Autor(es) (de 1 a 8)
\author{%
  Primeiro{(a)}~M.~Autor{(a)}%
  \AuthorInfo{%
    Email  = {author1@domain},%% (*)
    Lattes = {0000000000000001},%% (*)
    ORCID  = {0000-0000-0000-0001},%% (*) (CHKTEX 8)
    Inst   = {1},%
  }%
  \Inst{\faTalking}%% Palestrante (*)
  \and Segundo{(a)}~M.~Autor{(a)}%
  \AuthorInfo{%
    Email  = {author2@domain},%% (*)
    Lattes = {0000000000000002},%% (*)
    ORCID  = {0000-0000-0000-0002},%% (*) (CHKTEX 8)
    Inst   = {2},%
  }%
  \and Terceiro{(a)}~M.~Autor{(a)}%
  \AuthorInfo{%
    Email  = {author3@domain},%% (*)
    Lattes = {0000000000000003},%% (*)
    ORCID  = {0000-0000-0000-0003},%% (*) (CHKTEX 8)
    Inst   = {1},%
  }%
  \and\\ Quarto{(a)}~M.~Autor{(a)}%
  \AuthorInfo{%
    Email  = {author4@domain},%% (*)
    Lattes = {0000000000000004},%% (*)
    ORCID  = {0000-0000-0000-0004},%% (*) (CHKTEX 8)
    Inst   = {2},%
  }%
  \and Quinto{(a)}~M.~Autor{(a)}%
  \AuthorInfo{%
    Email  = {author5@domain},%% (*)
    Lattes = {0000000000000005},%% (*)
    ORCID  = {0000-0000-0000-0005},%% (*) (CHKTEX 8)
    Inst   = {1},%
  }%
}
%%%% Instituição(ões) e outros dados
\institute{%
  \Affil[1]{\UTFPRName, Cidade, Paraná, Brasil}%
  \and\Affil[2]{Instituição Parceira, Cidade, Estado, País}%
  \and\InfoList{\faEnvelope}{\MakeInfoList{Email}}%% (*)
  \and\InfoList{\faLattes}{\MakeInfoList{Lattes}}%% (*)
  \and\InfoList{\faOrcid}{\MakeInfoList{ORCID}}%% (*)
}
%%%% Data (atual por padrão) (*)
\date{}
%%%% Identificação do pôster (*)
\posterid{ID\@: EVNT\YearNum{-}0001}
%%%% Atribuição de Licença CC (*): {Ano}; {Autoria}
\CCLicenseAttribution{\YearNum}{P. M. Autor{(a)} \textit{et al.}}

%% Arquivo(s) de referências
\addbibresource{utfpr-poster.bib}

%% Início do documento
\begin{document}

\begin{frame}%% Início do quadro (pôster) >>>

\begin{columns}%% Cabeçalho (HDR) >>>
\dimen1 = 188mm%
\begin{column}{0.13\textwidth}%% HDR-C1 >>>
\raggedright%
\begin{minipage}[t][\dimen1]{\linewidth}
\vspace*{0pt}
\Logo{Evento}%% Logomarca evento
\vfill%
\Logo{Org. Promotora}%% Logomarca organização promotora
\end{minipage}
\end{column}%% <<< HDR-C1
\hfill%
\begin{column}{0.7\textwidth}%% HDR-C2 >>>
\titlepage%
\end{column}%% <<< HDR-C2
\hfill%
\begin{column}{0.13\textwidth}%% HDR-C3 >>>
\raggedleft%
\begin{minipage}[t][\dimen1]{\linewidth}
\vspace*{0pt}
\Logo{utfpr-d}%% Logomarca UTFPR ou campus (fundo claro ou escuro)
\vfill%
\Logo{Inst. Parceira}%% Logomarca instituição parceira
\end{minipage}
\end{column}%% <<< HDR-C3
\end{columns}%% <<< HDR

\vspace*{\stretch{1000}}

\begin{columns}%% L1 >>>
\begin{column}{0.49\textwidth}%% L1-C1 >>>
\begin{block}{Introdução}%% L1-C1-B1 >>>
Este pôster foi desenvolvido no modelo \UTFPR-Poster, baseado na classe \LaTeX\ \href{https://www.ctan.org/pkg/beamer/}{Beamer\LinkIcon} e usando o pacote \LaTeX\ \href{https://ctan.org/pkg/beamerposter/}{\texttt{beamerposter}\LinkIcon}.\par%
Exemplo de lista de itens numerados:
\begin{enumerate}
\item item numerado 1;
\item item numerado 2;
\item item numerado 3.
\end{enumerate}
\begin{alertblock}{\faInfoCircle\ Informações e dicas sobre \TeX/\LaTeX}%% L1-C1-B1-B1 >>>
\begin{itemize}
\selectlanguage{english}
\item \href{https://www.latex-project.org/}{\LaTeX\ Project\LinkIcon}.
\item \href{https://www.ctan.org/}{Comprehensive \TeX\ Archive Network (CTAN)\LinkIcon}.
\item \href{https://www.tug.org/}{\TeX\ Users Group (TUG)\LinkIcon}.
\item \href{https://en.wikibooks.org/wiki/LaTeX/}{\LaTeX\ \textemdash\ Wikibooks\LinkIcon}.
\item \href{https://tex.stackexchange.com/}{\TeX-\LaTeX\ Stack Exchange\LinkIcon}.
\end{itemize}
\end{alertblock}%% <<< L1-C1-B1-B1
\end{block}%% <<< L1-C1-B1
\begin{block}{Revisão da Literatura}%% L1-C1-B2 >>>
Exemplos de referências podem ser observados nas citações indiretas:
\begin{itemize}
\item Implícita: \ldots\ \cite{Ekenstein1997}.
\item Explícita: \textcite{Wizentier1992} analisaram\ldots%
\end{itemize}
Citações e referências podem ser inseridas neste documento usando os comandos do pacote \href{https://ctan.org/pkg/biblatex/}{Bib\LaTeX\LinkIcon}, conforme exemplos no arquivo-fonte deste modelo.
Os dados de cada referência podem ser obtidos de um arquivo \href{https://www.bibtex.org/}{Bib\TeX\LinkIcon} (\texttt{*.bib}), geralmente na própria página de acesso ou \ENLang{download} da publicação (artigos, livros, etc.) ou, ainda, a partir do Google Acadêmico, etc.
\begin{alertblock}{\faInfoCircle\ Ferramentas para gerar ou editar entradas \href{https://www.bibtex.org/}{Bib\TeX\LinkIcon}}%% L1-C1-B2-B1 >>>
\begin{itemize}
\selectlanguage{english}
\item[\scriptsize\faTools] \href{https://zbib.org/}{ZoteroBib\LinkIcon}.
\item[\scriptsize\faTools] \href{https://truben.no/latex/bibtex/}{Bib\TeX\ Editor\LinkIcon}.
\end{itemize}
\end{alertblock}%% <<< L1-C1-B2-B1
\end{block}%% <<< L1-C1-B2
\end{column}%% <<< L1-C1
\hfill%
\begin{column}{0.49\textwidth}%% L1-C2 >>>
\begin{block}{Material e Métodos}%% L1-C2-B1 >>>
Uma equação como $y = a x^2 + b x + c$ pode ser inserida ao longo do texto de um parágrafo usando o ambiente \LaTeX\ \texttt{math} (ou o atalho \LaTeX\ \texttt{\textbackslash(\ldots\textbackslash)} ou o atalho \TeX\ \texttt{\$\ldots\$}).
Por outro lado, a seguinte equação (não numerada) pode ser inserida em uma linha própria usando o ambiente \LaTeX\ \texttt{displaymath} (ou o atalho \LaTeX\ \texttt{\textbackslash[\ldots\textbackslash]}):
\[
\frac{\mathrm{d} y}{\mathrm{d} x} = \gamma \operatorname{sen} x
\]
A Equação~\eqref{eq:fx} foi inserida usando o ambiente \LaTeX\ \texttt{equation} e numerada automaticamente:
\begin{equation}%
\label{eq:fx}
f{\left(x\right)} = \frac{1}{\alpha} \int_0^L \left(\frac{x^2}{2} - \frac{x^3}{3}\right) \mathrm{d} x
\end{equation}
\begin{alertblock}{\faInfoCircle\ Ferramentas para gerar ou editar equações em \LaTeX}%% L1-C2-B1-B1 >>>
\begin{itemize}
\selectlanguage{english}
\item[\scriptsize\faTools] \href{https://formulasheet.com/}{Formula Sheet\LinkIcon}.
\item[\scriptsize\faTools] \href{https://www.tutorialspoint.com/latex_equation_editor.htm}{\LaTeX\ Equation Editor (\textit{by} Tutorials Point)\LinkIcon}.
\end{itemize}
\end{alertblock}%% <<< L1-C2-B1-B1
\end{block}%% <<< L1-C2-B1
\begin{block}{Resultados e Discussão}%% L1-C2-B2 >>>
A Figura~\ref{fig:campi-map}\footnote{Possui um código QR contendo um URL.} apresenta um mapa com a localização dos campi da UTFPR\@.
\begin{figure}[!htb]
\caption{Localização dos campi da UTFPR}%
\label{fig:campi-map}
\savebox0{\includegraphics[height = 100mm]{fig-campi-map}}
\usebox0%
\llap{%
  \raisebox{\ht0-\height}{%
    \href{https://www.utfpr.edu.br/campus}{%
      \includegraphics[height = 20mm]{fig-campi-map-qr-code}%
    }%
  }%
}
\SourceOrNote{\textcite{UTFPR2017}}
\end{figure}
\end{block}%% <<< L1-C2-B2
\end{column}%% <<< L1-C2
\end{columns}%% <<< L1

\begin{columns}%% L2 >>>
\begin{column}{\textwidth}%% L2-C1 >>>
\begin{block}{Resultados e Discussão (Continuação)}%% L2-C1-B1 >>>
\begin{column}[t]{0.3\linewidth}%% L2-C1-B1-C1 >>>
A Figura~\ref{fig:pg-campus}\footnote{Possui um código QR contendo um URL.} apresenta uma fotografia da fachada do campus Ponta Grossa da UTFPR\@.
\begin{figure}[!htb]
\caption{Fachada do campus Ponta Grossa da UTFPR}%
\label{fig:pg-campus}
\savebox0{\includegraphics[width = 0.75\linewidth]{phot-pg-campus}}
\usebox0%
\llap{%
  \raisebox{\ht0-\height}{%
    \href{https://www.utfpr.edu.br/campus/pontagrossa}{%
      \includegraphics[height = 20mm]{phot-pg-campus-qr-code}%
    }%
  }%
}
\SourceOrNote{\textcite{UTFPR2018}}
\end{figure}
\end{column}%% <<< L2-C1-B1-C1
\YellowVRule\hfill%
\begin{column}[t]{0.3\linewidth}%% L2-C1-B1-C2 >>>
A Figura~\ref{fig:t-x}\footnote{Gráfico produzido no ambiente \LaTeX\ \texttt{tikzpicture} do pacote \LaTeX\ \texttt{tikz} a partir do arquivo \texttt{grph-t-x.tex} em \texttt{./Figures/}.} foi inserida usando o ambiente \LaTeX\ \texttt{figure} e numerada automaticamente.
\begin{figure}[!htb]
\caption{Exemplo de legenda de figura}%
\label{fig:t-x}
\includegraphics[height = 140mm]{grph-t-x}
\SourceOrNote{autoria própria (\YearNum)}
\end{figure}
\end{column}%% <<< L2-C1-B1-C2
\YellowVRule\hfill%
\begin{column}[t]{0.375\linewidth}%% L2-C1-B1-C3 >>>
A Tabela~\ref{tab:L-dimensions} foi inserida usando o ambiente \LaTeX\ \texttt{table} e numerada automaticamente.
\begin{table}[!htb]
\caption{Exemplo de legenda de tabela}%
\label{tab:L-dimensions}
\begin{tabularx}{\linewidth}{@{}XYYYY@{}}
\toprule%
\multicolumn{1}{Y}{\textbf{Caso}} &
\MathBF{L}                        &
\MathBF{L^2}                      &
\MathBF{L^3}                      &
\multicolumn{1}{Y}{\MathBF{L^4}}  \\
\multicolumn{1}{Y}{}                      &
\MathBF{(\Unit{m})}                       &
\MathBF{(\Unit{m^2})}                     &
\MathBF{(\Unit{m^3})}                     &
\multicolumn{1}{Y}{\MathBF{(\Unit{m^4})}} \\ \midrule%
A & 1 & 1  & 1   & 1   \\
B & 2 & 4  & 8   & 16  \\
C & 3 & 9  & 27  & 81  \\
D & 4 & 16 & 64  & 256 \\
E & 5 & 25 & 125 & 625 \\ \bottomrule%
\addlinespace[\abovecaptionskip]
\end{tabularx}
\SourceOrNote{autoria própria (\YearNum)}
\end{table}
\begin{alertblock}{\faInfoCircle\ Ferramentas para gerar ou editar tabelas em \LaTeX}%% L2-C1-B1-C3-B1 >>>
\begin{itemize}
\selectlanguage{english}
\item[\scriptsize\faTools] \href{https://www.tablesgenerator.com/}{Tables Generator\LinkIcon}.
\item[\scriptsize\faTools] \href{https://www.latex-tables.com/}{\LaTeX\ Tables Editor\LinkIcon}.
\end{itemize}
\end{alertblock}%% <<< L2-C1-B1-C3-B1
\end{column}%% <<< L2-C1-B1-C3
\end{block}%% <<< L2-C1-B1
\end{column}%% <<< L2-C1
\end{columns}%% <<< L2

\begin{columns}%% L3 >>>
\begin{column}{0.29\textwidth}%% L3-C1 >>>
\begin{block}{Conclusões}%% L3-C1-B1 >>>
As conclusões ou considerações finais podem ser apresentadas como uma lista de itens, enfatizando as contribuições do trabalho:
\begin{itemize}
\item Primeiro item de conclusão.
\item Segundo item de conclusão.
\item Terceiro item de conclusão.
\end{itemize}
\end{block}%% <<< L3-C1-B1
\end{column}%% <<< L3-C1
\hfill%
\begin{column}{0.69\textwidth}%% L3-C2 >>>
\begin{block}{Referências}%% L3-C2-B1 >>>
\printbibliography[heading = none]
\end{block}%% <<< L3-C2-B1
\end{column}%% <<< L3-C2
\end{columns}%% <<< L3

\begin{columns}%% L4 >>>
\dimen1 = 30mm%
\begin{column}{0.6\textwidth}%% L4-C1 >>>
\begin{block}{Agradecimentos}%% L4-C1-B1 >>>
\footnotesize%
\bigskip%
\Logo*{capes-l}%% Logomarca CAPES
\hfill%
\Logo*{cnpq-l}%% Logomarca CNPq
\hfill%
\Logo*{fa-l}%% Logomarca Fundação Araucária
\hfill%
\Logo*{Extra}%% Logomarca extra
\hfill%
\Logo*{public-education-l}%% Logomarca extra (Ensino Público)
\hspace*{10mm}%
\Logo*{utfpr-l}%% Logomarca UTFPR ou campus (fundo claro ou escuro)
\bigskip%
\end{block}%% <<< L4-C1-B1
\end{column}%% <<< L4-C1
\hfill%
\begin{column}{0.3\textwidth}%% L4-C2 >>>
\begin{exampleblock}{\faTalking\ Palestrante}%% L4-C2-B1 >>>
\footnotesize%
\bigskip%
\begin{minipage}{0.75\dimen1}
\color{UTFPRYellow}
\setkeys{draftfigure}{content = {\faUserTie\\Inserir\\Foto}, size = {tiny}}
\includegraphics[width = 0.75\dimen1, height = \dimen1, draft]{example-image}
\end{minipage}%
\hfill%
\begin{minipage}{\linewidth-1em-0.75\dimen1}
\InfoList{\faUserTie}{Nome Completo do{(a)} Autor{(a)}}\par%
\InfoList{\faUniversity}{\UTFPRName}\par%
\InfoList{\faEnvelope}{\email{author1@domain}}
\end{minipage}%
\bigskip%
\end{exampleblock}%% <<< L4-C2-B1
\end{column}%% <<< L4-C2
\hfill%
\begin{column}{0.06\textwidth}%% L4-C3 >>>
\begin{alertblock}{\faPlus\ \faInfoCircle}%% L4-C3-B1 >>>
\footnotesize%
\centering%
\bigskip%
\href{https://www.utfpr.edu.br/}{\includegraphics[height = \dimen1]{utfpr-qr-code}}%
\bigskip%
\end{alertblock}%% <<< L4-C3-B1
\end{column}%% <<< L4-C3
\end{columns}%% <<< L4

\vspace*{\stretch{1000}}

\null%

\begin{columns}%% Rodapé (FTR) >>>
\begin{column}{0.49\textwidth}%% FTR-C1 >>>
\RespNotice[Declaração de Responsabilidade]{o{(s)} autor{(es)} é{(são)} o{(s)} único{(s)} responsável{(eis)} pelas informações contidas neste documento.}
\end{column}%% <<< FTR-C1
\hfill%
\begin{column}{0.49\textwidth}%% FTR-C2 >>>
\CCLicence%
\end{column}%% <<< FTR-C2
\end{columns}%% <<< FTR

\null%

\end{frame}%% <<< Fim do quadro (pôster)

%% Fim do documento
\end{document}
